%%%%%%%%%%%%%%%%%%%%%%%%%%%%%%%%%%%%%%%%%
% 陈嘉显 - AI大模型算法开发实习生简历
% 基于 resume.cls 模板（终极稳定版）
%%%%%%%%%%%%%%%%%%%%%%%%%%%%%%%%%%%%%%%%%

\documentclass{resume} 
\usepackage{ctex}
\usepackage[left=0.4in,top=0.2in,right=0.4in,bottom=0.2in]{geometry}
\usepackage{amssymb}

\name{陈嘉显}
% 使用 \makebox[\textwidth][l] 强制左对齐且不换行，彻底解决 address 居中换行问题
\address{\makebox[\textwidth][l]{现居住地：北京市 \hfill 13959591670 \hfill 1251828394@qq.com}}

\begin{document}

%----------------------------------------------------------------------------------------
%	教育背景
%----------------------------------------------------------------------------------------

\begin{rSection}{Education}
\noindent\textbf{北京邮电大学} \hfill \textbf{网络空间安全（硕士）} \hfill 2025年7月 -- 至今 \\
\noindent\textbf{福州大学} \hfill \textbf{信息安全（本科）} \hfill 2020年9月 -- 2024年6月 \\
GPA: 3.15/4.00（专业前25\%），获三次校奖学金（2020--2021 第一/二学期、2021--2022 第一学期）
\end{rSection}

%----------------------------------------------------------------------------------------
%	项目经历
%----------------------------------------------------------------------------------------

\begin{rSection}{Projects}
\begin{rSubsection}{NanoChat-Lab —— 轻量化 RAG \& Agent 系统（含 LoRA 微调）}{2026年3月}{个人项目}{https://github.com/xmfcjx/NanoChat-Lab}
\item \textbf{项目介绍}：在 4GB 显存（GTX 1650）约束下，基于 Qwen2.5-1.5B 构建端到端的 RAG + Agent + Memory 系统，并通过 LoRA 微调 + Prompt 压缩实现”更准、更快、更省”的工具调用能力，全部模块自研实现，不依赖 LangChain。
\item \textbf{LoRA 微调与 Prompt 压缩}：基于 3000 条 SFT 数据在云端完成 LoRA 微调（rank=8, 300 steps），配合自研的精简 System Prompt（149 tokens vs 原始 282 tokens），在 200 条 SBS 对比评估中实现工具准确率 +5\%、答案准确率 +8.5\%、格式正确率 +8\%，同时 Token 节省 47\%、延迟降低 1.25s。通过 verify 模式诊断确认性能瓶颈来自 Prompt 设计而非微调质量。
\item \textbf{自研混合检索架构}：基于 NumPy 自研轻量向量库，构建”Vector + BM25 混合召回 + Cross-Encoder 重排”双路检索体系，消融实验中 Top-3 命中率达 92\%，有效解决单一向量检索在长文本及专业关键词场景下的性能短板。
\item \textbf{InputClassifier 路由 + Agent 直通设计}：通过正则匹配与语义识别实现任务预分类，设计 Agent 直通模式跳过 ReAct 循环；大规模测试表明直通模式将准确率提升约 40\%，延迟降低近 800 倍，有效解决小模型 ReAct 循环不稳定性问题。
\item \textbf{多维评测体系}：构建 SBS 对比评估框架，支持工具准确率/答案准确率/格式正确率/延迟四维度自动评测，内置 DPO 偏好数据生成与 verify 诊断模式，实现”训练-评估-诊断-迭代”闭环。
\item \textbf{模型量化与部署}：支持 fp16/int8/int4/fp4/GGUF 五种量化模式，int4 量化下显存占用仅 2.8GB，配合 LoRA adapter 热加载实现低资源设备上的高效推理。
\end{rSubsection}
\end{rSection}

%----------------------------------------------------------------------------------------
%	竞赛获奖
%----------------------------------------------------------------------------------------

\begin{rSection}{Awards \& Competitions}
{\bf 高教社杯全国大学生数学建模竞赛 —— 全国二等奖} \hfill {2022年9月 -- 11月}
\begin{itemize}
    \item 论文《基于贪心算法的无人机纯方位无源定位技术》获全国二等奖，并发表于 EI 会议：\textit{Proc. IEEE ACDP}, 2023. 通过贪心算法和迭代法将误差函数控制在 $10^{-25}$ 数量级，实现有限次数内的无源精确定位。
\end{itemize}
% \vspace{-0.8em}
\textbf{团队贡献}：团队任务分配；对论文模型构建提供数理证明。
\end{rSection}

%----------------------------------------------------------------------------------------
%	毕业论文
%----------------------------------------------------------------------------------------

\begin{rSection}{Thesis}
{\bf 毕业论文-通过捕获 EMR 实现个性化临床特征嵌入研究} 
\begin{itemize}
    \item 针对现有算法模型在个性化治疗中的不足，提出一种优化框架。该框架引入多通道 GRU 分别学习动态特征表示，并结合特征编码器整合静态信息与动态记录，同时通过 Multi-Head Attention 与协方\textbf{}差最小化策略增强特征表示的多样性，最终有效提升了个性化预测的准确性和泛化能力。
\end{itemize}
\end{rSection}

%----------------------------------------------------------------------------------------
%	专业技能
%----------------------------------------------------------------------------------------

\begin{rSection}{Skills}
\textbf{大模型与深度学习:} 深入理解 Transformer 架构；熟悉 PyTorch 框架及底层架构，具备 Llama 类模型从零训练经验；熟悉 LoRA/QLoRA 微调流程与 SFT/DPO 对齐方法；熟悉 ReAct 智能体逻辑与多工具调用；具备 CS336 大模型工程实践经验。 \\
\textbf{RAG 与搜索优化:} 熟悉 RAG 全链路开发；熟悉向量检索与 BM25 混合召回策略；熟练应用 Cross-Encoder 进行 Rerank；具备基于 NumPy 手写轻量化向量数据库与动态分块算法的底层实现能力。 \\
\textbf{编程与工具:} 熟悉 Python 与 C/C++ 编程，具备扎实的算法与数据结构基础；熟悉 HuggingFace Transformers/PEFT/bitsandbytes 等模型训练与推理工具链。
\textbf{     }

\end{rSection}

\end{document}
